\documentclass[11pt,a4paper]{article}

\usepackage[T1]{fontenc}
\usepackage[utf8]{inputenc}
\usepackage{lmodern}
\usepackage{amssymb,amsmath,amsthm,mathtools}
\usepackage[margin=2.5cm]{geometry}
\usepackage{parskip}
\usepackage[hidelinks]{hyperref}
\usepackage{microtype}

% Theorem environments (arXiv / journal style)
\theoremstyle{plain}
\newtheorem{theorem}{Theorem}[section]
\newtheorem{proposition}[theorem]{Proposition}
\newtheorem{lemma}[theorem]{Lemma}
\newtheorem{corollary}[theorem]{Corollary}

\theoremstyle{definition}
\newtheorem{definition}[theorem]{Definition}
\newtheorem{example}[theorem]{Example}
\newtheorem{remark}[theorem]{Remark}

% Notation
\newcommand{\Shape}{\mathrm{Shape}}
\newcommand{\Pos}{\mathrm{Pos}}
\newcommand{\Type}{\mathbf{Type}}
\newcommand{\id}{\mathrm{id}}
\newcommand{\src}{\mathrm{src}}
\newcommand{\tgt}{\mathrm{tgt}}
\newcommand{\rootpos}{\mathrm{root}}
\newcommand{\sub}{\mathrm{sub}}
\newcommand{\shift}{\mathrm{shift}}

\title{Aggregation as a Routing Cofunctor}
\author{Extracted from a machine-checked development}
\date{}

\begin{document}
\maketitle

\begin{abstract}
We construct a routing policy from the concrete eleven-facility Sawasdee
Cellars wine network to a two-node aggregate network (Australia to Thailand)
and verify that it is a morphism of directed containers. The shape map sends
each facility to its geographic region; the position map lifts each aggregate
route to a canonical concrete route out of the given facility. Shape
compatibility, root preservation, and shift compatibility are established by
finite case analysis. The proofs of all stated results have been
machine-checked; we record only the mathematical content.
\end{abstract}

\tableofcontents
\bigskip

%======================================================================
\section{Introduction}
%======================================================================

Aggregation in logistics replaces a fine-grained network of facilities by a
coarse network of regions, while retaining enough route structure that
policies expressed at the regional level can be interpreted concretely. In the
language of directed containers, such a policy is a morphism from the concrete
network to the aggregate network: a map on shapes (facilities to regions)
together with a \emph{lift} of aggregate positions to concrete positions,
compatible with roots, sub-shapes, and shifts.

The present note isolates one fully explicit instance. Starting from the
Sawasdee Cellars wine network $\mathcal{N}_{\mathrm{wine}}$ (eleven facilities,
ten shipping legs), we form a two-node aggregate digraph $G_{\mathrm{agg}}$
with regions $\{\mathrm{australia},\, \mathrm{thailand}\}$ and a single
aggregate leg. A routing policy
\[
  \varphi \;:\;
  \mathcal{N}_{\mathrm{wine}}
  \;\Longrightarrow\;
  \mathrm{network}(G_{\mathrm{agg}})
\]
is defined by sending each facility to its region and lifting the aggregate
stay-put and Australia-to-Thailand routes to canonical concrete routes. The
three morphism laws are checked by exhaustive case analysis on the finite set
of facilities.

The development works throughout in a single universe of types. Proofs are
indicated by the phrase ``by lean,'' meaning they have been checked by a proof
assistant rather than written out in prose.

\paragraph{What this note claims.}
In ordinary mathematical language: coarsening the wine network to two
geographic regions extends to a directed-container morphism. Regional route
data lift canonically to facility-level routes, and the lift preserves the
algebraic structure of roots and path concatenation.

%======================================================================
\section{Directed-container morphisms}
%======================================================================

We recall the data of a morphism of directed containers in the form used
below. Let $\mathcal{C}$ and $\mathcal{D}$ be directed containers, with shape
types $S_{\mathcal{C}}$ and $S_{\mathcal{D}}$, position families $P_{\mathcal{C}}$
and $P_{\mathcal{D}}$, and structure maps $\rootpos$, $\sub$, $\shift$.

\begin{definition}[Directed-container morphism]\label{def:dc-morphism}
A \textbf{morphism of directed containers}
$\varphi : \mathcal{C} \Rightarrow \mathcal{D}$ consists of
\begin{enumerate}
  \item a map on shapes $u : S_{\mathcal{C}} \to S_{\mathcal{D}}$, and
  \item for each shape $s \in S_{\mathcal{C}}$, a map on positions
        \[
          \ell_s \;:\;
          P_{\mathcal{D}}\bigl(u(s)\bigr)
          \;\longrightarrow\;
          P_{\mathcal{C}}(s)
        \]
        (a \emph{lift} of $\mathcal{D}$-positions to $\mathcal{C}$-positions),
\end{enumerate}
subject to three laws:
\begin{enumerate}
  \item[\textbf{(M0)}] \emph{Shape compatibility.} For every $s$ and every
        position $\alpha \in P_{\mathcal{D}}\bigl(u(s)\bigr)$,
        \[
          u\Bigl(\sub_{\mathcal{C}}\bigl(s,\, \ell_s(\alpha)\bigr)\Bigr)
          \;=\;
          \sub_{\mathcal{D}}\bigl(u(s),\, \alpha\bigr).
        \]
  \item[\textbf{(M1)}] \emph{Root preservation.} For every~$s$,
        \[
          \ell_s\bigl(\rootpos_{\mathcal{D}}(u(s))\bigr)
          \;=\;
          \rootpos_{\mathcal{C}}(s).
        \]
  \item[\textbf{(M2)}] \emph{Shift compatibility.} For every~$s$, every
        $\alpha \in P_{\mathcal{D}}\bigl(u(s)\bigr)$, and every
        $\kappa \in P_{\mathcal{D}}\bigl(
          \sub_{\mathcal{D}}(u(s),\alpha)\bigr)$,
        \[
          \ell_s\bigl(
            \shift_{\mathcal{D}}\bigl(u(s),\, \alpha,\, \kappa\bigr)
          \bigr)
          \;=\;
          \shift_{\mathcal{C}}\Bigl(
            s,\;
            \ell_s(\alpha),\;
            \ell_{\sub_{\mathcal{C}}(s,\ell_s(\alpha))}
              \bigl(\kappa'\bigr)
          \Bigr),
        \]
        where $\kappa'$ is $\kappa$ transported along the equality~(M0).
\end{enumerate}
\end{definition}

\begin{remark}
The variance on positions is covariant for directed-container morphisms in this
formulation (lifts of target positions to source positions), matching the
cofunctorial reading of aggregation: coarse routes are refined to fine routes.
\end{remark}

%======================================================================
\section{The aggregate network}
%======================================================================

\begin{definition}[Aggregate digraph]\label{def:gagg}
Let $G_{\mathrm{agg}}$ be the directed graph with
\begin{enumerate}
  \item nodes $\{\mathrm{australia},\, \mathrm{thailand}\}$, and
  \item a single edge $\mathrm{au2th} : \mathrm{australia} \to
        \mathrm{thailand}$.
\end{enumerate}
Write $\mathcal{N}_{\mathrm{agg}} := \mathrm{network}(G_{\mathrm{agg}})$ for the
associated directed container.
\end{definition}

\begin{remark}[Aggregate positions]
The positions of $\mathcal{N}_{\mathrm{agg}}$ are elementary:
\begin{itemize}
  \item at $\mathrm{australia}$: the empty route
        $\langle\mathrm{australia},\, \mathrm{nil}\rangle$ and the single-leg
        route $\langle\mathrm{thailand},\, \mathrm{cons}(\mathrm{au2th},\,
        \mathrm{nil})\rangle$;
  \item at $\mathrm{thailand}$: only the empty route
        $\langle\mathrm{thailand},\, \mathrm{nil}\rangle$.
\end{itemize}
\end{remark}

%======================================================================
\section{The concrete wine network and the region map}
%======================================================================

Recall the Sawasdee Cellars network $\mathcal{N}_{\mathrm{wine}} =
\mathrm{network}(G_{\mathrm{wine}})$: eleven facilities and ten shipping legs
from vineyards in South Australia through an export agency and sea freight to
Thai distributors and retail shops.

\begin{definition}[Region assignment]\label{def:region}
Define a map $\mathrm{region} : \mathrm{Facility} \to
\{\mathrm{australia},\, \mathrm{thailand}\}$ by
\begin{align*}
  \mathrm{region}(\mathrm{vineyard}_i)
  &\;=\; \mathrm{australia}
    &&(i=1,2,3), \\
  \mathrm{region}(\mathrm{sawasdee})
  &\;=\; \mathrm{australia}, \\
  \mathrm{region}(\mathrm{melbournePort})
  &\;=\; \mathrm{australia}, \\
  \mathrm{region}(\mathrm{thaiPort})
  &\;=\; \mathrm{thailand}, \\
  \mathrm{region}(\mathrm{distributor}_j)
  &\;=\; \mathrm{thailand}
    &&(j=1,2), \\
  \mathrm{region}(\mathrm{retail}_k)
  &\;=\; \mathrm{thailand}
    &&(k=1,2,3).
\end{align*}
\end{definition}

Thus the five Australian-side facilities map to $\mathrm{australia}$ and the
six Thai-side facilities map to $\mathrm{thailand}$. The sea leg of the
concrete network is the unique edge that crosses the regional boundary.

%======================================================================
\section{Lifting aggregate routes}
%======================================================================

Write $\mathrm{AgPos}(r)$ for $\Pos(\mathcal{N}_{\mathrm{agg}})(r)$ and
$\mathrm{Route}(f)$ for $\Pos(\mathcal{N}_{\mathrm{wine}})(f)$.

\begin{definition}[Position lift]\label{def:lift}
Define, for each facility~$f$, a map
\begin{equation}
  \ell_f \;:\;
  \mathrm{AgPos}\bigl(\mathrm{region}(f)\bigr)
  \;\longrightarrow\;
  \mathrm{Route}(f)
\end{equation}
by the following canonical choices.
\begin{enumerate}
  \item \textbf{Australian facilities.} If $f$ is Australian and the aggregate
        route is the empty route at $\mathrm{australia}$, lift to the empty
        concrete route at~$f$. If the aggregate route ends in
        $\mathrm{thailand}$, lift to the unique canonical concrete route from
        $f$ to $\mathrm{thaiPort}$ along the spine
        \[
          f \;\to\; \cdots \;\to\;
          \mathrm{melbournePort} \;\xrightarrow{\mathrm{sea}}\;
          \mathrm{thaiPort}
        \]
        (prefixing the legs $v_is$, $smp$, and/or $\mathrm{sea}$ as required by
        the facility). Concretely:
        \begin{itemize}
          \item from $\mathrm{vineyard}_i$: legs $v_is$, $smp$, $\mathrm{sea}$;
          \item from $\mathrm{sawasdee}$: legs $smp$, $\mathrm{sea}$;
          \item from $\mathrm{melbournePort}$: leg $\mathrm{sea}$.
        \end{itemize}
  \item \textbf{Thai facilities.} If $f$ is Thai, the only aggregate position
        is the empty route at $\mathrm{thailand}$; lift it to the empty
        concrete route at~$f$.
\end{enumerate}
\end{definition}

\begin{remark}
Because $G_{\mathrm{wine}}$ is a polytree along the Australia-to-Thailand spine
used here, the ``canonical'' concrete route from each Australian facility to
Thai Port is unambiguous. The lift does not attempt to choose among downstream
distributor or retail alternatives; those live entirely within the Thai region
and are invisible to the two-node aggregate.
\end{remark}

%======================================================================
\section{The morphism laws}
%======================================================================

\begin{proposition}[Shape compatibility, M0]\label{prop:m0}
For every facility~$f$ and every aggregate position
$\alpha \in \mathrm{AgPos}\bigl(\mathrm{region}(f)\bigr)$,
\begin{equation}
  \mathrm{region}\Bigl(
    \sub_{\mathcal{N}_{\mathrm{wine}}}\bigl(f,\, \ell_f(\alpha)\bigr)
  \Bigr)
  \;=\;
  \sub_{\mathcal{N}_{\mathrm{agg}}}\bigl(\mathrm{region}(f),\, \alpha\bigr).
\end{equation}
\end{proposition}

\begin{proof}
By lean.
\end{proof}

\begin{proposition}[Root preservation, M1]\label{prop:m1}
For every facility~$f$,
\begin{equation}
  \ell_f\bigl(
    \rootpos_{\mathcal{N}_{\mathrm{agg}}}\bigl(\mathrm{region}(f)\bigr)
  \bigr)
  \;=\;
  \rootpos_{\mathcal{N}_{\mathrm{wine}}}(f).
\end{equation}
\end{proposition}

\begin{proof}
By lean.
\end{proof}

\begin{proposition}[Shift compatibility, M2]\label{prop:m2}
For every facility~$f$, every aggregate position
$\alpha \in \mathrm{AgPos}\bigl(\mathrm{region}(f)\bigr)$, and every aggregate
continuation
$\kappa \in \mathrm{AgPos}\bigl(
  \sub_{\mathcal{N}_{\mathrm{agg}}}(\mathrm{region}(f),\alpha)\bigr)$,
the lift of the aggregate shift equals the concrete shift of the lifted
positions (with the continuation transported along Proposition~\ref{prop:m0}).
\end{proposition}

\begin{proof}
By lean.
\end{proof}

\begin{remark}[Idea of the case analysis]
All three laws are proved by matching on the finite type of facilities and on
the (at most two) aggregate positions at the image region. Representative
situations:
\begin{itemize}
  \item Empty aggregate route at an Australian facility: sub-shape stays
        Australian; lifting the root is the concrete root; shifting by a
        further Australia-to-Thailand leg reproduces the canonical concrete
        crossing.
  \item Australia-to-Thailand aggregate route: sub-shape is Thailand; the
        concrete lift ends at Thai Port; further aggregate continuations are
        necessarily the Thai empty route, and concatenating the empty suffix
        leaves the concrete route unchanged.
  \item Thai facilities: only empty routes appear, and all structure maps are
        trivial.
\end{itemize}
\end{remark}

\begin{theorem}[Routing policy]\label{thm:routing-policy}
The pair $\bigl(\mathrm{region},\, \ell\bigr)$ is a morphism of directed
containers
\begin{equation}
  \varphi \;:\;
  \mathcal{N}_{\mathrm{wine}}
  \;\Longrightarrow\;
  \mathcal{N}_{\mathrm{agg}}.
\end{equation}
\end{theorem}

\begin{proof}
By lean.
\end{proof}

%======================================================================
\section{Discussion}
%======================================================================

\paragraph{What has been established.}
Geographic aggregation of the Sawasdee Cellars network to two regions extends
to a directed-container morphism. The construction supplies a fully checked
instance of aggregation-as-cofunctor: coarse route algebra is interpreted by
lifting to fine routes along a fixed regional policy.

\paragraph{On canonicity.}
The position lift chooses, for each Australian facility, a single spine route
to Thai Port. Alternative concrete realisations (e.g.\ different downstream
branching inside Thailand) are invisible at the two-node aggregate level; a
finer aggregate would require a richer lift.

\paragraph{Scope.}
The development treats one concrete pair of networks and does not develop a
general theory of aggregation morphisms for arbitrary partitions of nodes. It
also does not address the interaction of routing policies with compositional
trackings (e.g.\ whether monoid-valued data descends along~$\varphi$).

\paragraph{Machine checking.}
Every definition and theorem statement above corresponds to a declaration in a
machine-checked formalisation. The phrase ``by lean'' indicates that the proof
has been verified by a proof assistant.

%======================================================================
\section*{Acknowledgements}
%======================================================================

The mathematical content follows the aggregation construction in the
supply-chain development that motivated this note. The presentation is a prose
rendering of a formal development intended for human review of the theorem
statements.

\begin{thebibliography}{9}

\bibitem{abbott2005}
M.~Abbott, T.~Altenkirch, and N.~Ghani.
\newblock Containers: constructing strictly positive types.
\newblock \emph{Theoretical Computer Science}, 342(1):3--27, 2005.

\bibitem{ahman2014}
D.~Ahman and T.~Uustalu.
\newblock Coalgebraic update lenses.
\newblock In \emph{MFPS}, volume~308 of \emph{ENTCS}, pages 25--48, 2014.

\end{thebibliography}

\end{document}
